%!Mode:: "TeX:UTF-8"
%!TEX encoding = UTF-8 Unicode
%!TEX TS-program = xelatex
\documentclass{ctexart}
\usepackage{bbm}
\usepackage{tikz}
\usepackage{amsmath,amssymb,amsthm,color,mathrsfs}
\usepackage{hyperref}
\usepackage{cleveref}
\usepackage{enumitem,anysize}%
\usepackage{expl3}
\marginsize{1in}{1in}{1in}{1in}%
\theoremstyle{remark}
\newtheorem*{claim*}{Claim}
\newtheorem{lemma}{Lemma}

\DeclareMathOperator{\ot}{ordertype}
\DeclareMathOperator{\dom}{dom}
\DeclareMathOperator{\ran}{ran}
\DeclareMathOperator{\field}{field}
\DeclareMathOperator{\Ord}{Ord}

\newcommand{\cc}{\mathfrak{c}}

\def\email#1{{\texttt{#1}}}
\newcommand{\isep}[1][0pt]{\addtolength{\itemsep}{#1}}
\newcounter{problem}
\renewcommand{\theproblem}{\Roman{problem}}
\newenvironment{problem}{\refstepcounter{problem}\noindent\color{blue}\begin{tikzpicture}[baseline]%
\node at (-0.02em,0.3em) {$\mathbb{P}$};%
\node[scale=0.7] at (0.2em,-0.0em) {R};%
\node[scale=0.7] at (0.6em,0.4em) {O};%
\node[scale=0.8] at (1.05em,0.25em) {B};%
\node at (1.55em,0.3em) {L};%
\node[scale=0.7] at (1.75em,0.45em) {E};%
\node at (2.35em,0.3em) {M};%
\end{tikzpicture}\theproblem}{}
\crefname{problem}{Problem}{Problem}
%\renewcommand\theprob{{\Roman{problem}}}
\newcommand\mysolution{\begin{tikzpicture}[baseline]%
\node at (-0.04em,0.3em) {$\mathbb{S}$};%
\node[scale=0.7] at (0.35em,0.4em) {O};%
\node at (0.7em,0.3em) {\textit{L}};%
\node[scale=0.7] at (0.95em,0.4em) {U};%
\node[scale=1.1] at (1.19em,0.32em){T};%
\node[scale=0.85] at (1.4em,0.24em){I};%
\node at (1.9em,0.32em){$\mathcal{O}$};%
\node[scale=0.75] at (2.3em,0.21em){\texttt{N}};%
\end{tikzpicture}}
\newenvironment{solution}{\begin{proof}[\mysolution]}{\end{proof}}

%%%%%%%%
\newcommand{\calL}{\mathcal{L}}

\newcommand\<{\langle}
\renewcommand\>{\rangle}
\newcommand\eneg{\mathcal{E}_{\neg}}
\newcommand\eto{\mathcal{E}_{\to}}
\newcommand\N{\mathbb{N}}
\newcommand\subini{\subsetneqq_{init}}
\def\to{\rightarrow}
\newcommand{\calA}{\mathcal{A}}
\newcommand{\xor}{\vee}
\newcommand{\bor}{\bigvee}
\newcommand{\band}{\bigwedge}
\newcommand{\xand}{\wedge}
\newcommand{\minus}{\mathbin{\backslash}}
\newcommand{\calF}{\mathcal{F}}
\newcommand\mi[1]{\mathscr{P}(#1)}
\newcommand{\calB}{\mathcal{B}}
\newcommand{\heq}{\mathop{\hat{=}}}
\newcommand{\hneq}{\mathop{\hat{\neq}}}
\newcommand{\calR}{\mathcal{R}}
\newcommand{\hle}{\mathop{\hat{<}}}
\newcommand{\R}{\mathbb{R}}
\renewcommand{\C}{\mathbb{C}}
\newcommand{\calM}{\mathcal{M}}
\newcommand{\calN}{\mathcal{N}}
\newcommand{\frA}{\mathfrak{A}}
\newcommand{\card}{\mathrm{card}}
\newcommand{\frM}{\mathfrak{M}}
\newcommand{\Th}{\mathrm{Th}}
\newcommand{\ecd}{\mathrm{EC}_\Delta}
\newcommand{\Q}{\mathbb{Q}}
\newcommand{\Z}{\mathbb{Z}}
\newcommand{\oto}{\leftrightarrow}
\newcommand{\hin}{\hat{\in}}
\newcommand{\fun}[2]{{}^{#1}#2}
\newcommand{\A}{\mathbbm{A}}
\newcommand{\e}{\mathrm{e}}
\newcommand{\gl}{\mathrm{GL}}
\newcommand{\im}{\mathrm{Im}}
\newcommand{\re }{\mathrm{Re }}
\newcommand{\rank}{\mathrm{rank}}







\begin{document}
\title{$\mathbb{GROUP\ REPRESENTATION}$}
\author{白永乐\\ %% 姓名
SID: 202011150087\\ %% 学号
\email{202011150087"mail.bnu.edu.cn}} %% 电邮
\maketitle
\begin{problem}\label{pro:1}
    Find an $n-$dimensional real matrix representation of $(\R,+)$. 
\end{problem}

\begin{solution}
    Let:
    \begin{equation}
        \phi : \R\to \gl_n(\R),x\mapsto e^x I_n
    \end{equation}
    Then obviously $(M_n(\R),\phi)$ is an $n-$dimensional real matrix representation of $(\R,+)$. 
\end{solution}

\begin{problem}\label{pro:2}
    Find an infinitely dimensional representation of $(\R,+)$. 
\end{problem}
\begin{solution}
    Let:
    \begin{equation}
        \phi:\R\to\gl(\R[x]),\phi(a)(f)(x):=f(a+x)
    \end{equation}
    Then it's easy to know $(\phi,\R[x])$ is a representation of $(\R,+)$. 
\end{solution}

\begin{problem}
    Determine whether the representation is faithful or not in 
    Example 1~4 and 
    \Cref{pro:1}, 
    \Cref{pro:2}. 
\end{problem}

\begin{solution}
    \begin{itemize}
        \item Example 1: When $a=0$ we have $f_0(x)=\e^0=1$, 
        so it's obviously not faithful. When $a\neq 0$, we can easily get $f_a$ is injective, so it's faithful. 
        \item Example 2: When $a=0$ we have $f_0(x)=\e^0=1$, 
        so it's obviously not faithful. When $a\in \R^*$, let $x=y+\frac{2\pi}{a}$, then $f_a(x)=f_a(y)$, so $f_a$ is not faithful. When $a\in\C\minus \R$, we get $|f_a(x)|=\e^{-\im(a) x}$ is injective, so $f_a$ is injective, thus faithful. 
        \item Example 3: Obviously $\phi(x+2\pi )=\phi(x)$ so it's not faithful. 
        \item Example 4: Consider $f\in\R_n[x],f(x)=x$, then for $a\neq b$, we have $\phi(a)(f)=x+a\neq x+b=\phi(b)(f)$, so $\phi(a)\neq \phi(b)$, thus $\phi$ is faithful. 
        \item \Cref{pro:1}: Obviously it's faithful.
        \item \Cref{pro:2}: Obviously it's faithful. 
    \end{itemize}
\end{solution}

\begin{problem}
    Let $\lambda\in \C^*$, and:
    \begin{equation}
        \phi_{\lambda}:(\Z,+)\to\C^*,n\mapsto \lambda^n
    \end{equation}
    Prove that $\phi_\lambda$ is an $1-$dimensional complex representation of $(\Z,+)$, and find when it's faithful.   
\end{problem}

\begin{solution}
    For $m,n\in \Z$, we have $\phi_{\lambda}(m+n)=\lambda^{m+n}=\lambda^m \lambda^n =\phi_{\lambda}(m)\phi_{\lambda}(n)$, so $\phi_{\lambda}$ is complex representation. Obviously $\C^*\cong \gl(\C)$, so it has dimension one. 

    Noting $\phi_\lambda$ is faithful$\iff \ker(\phi)=\{0\}\iff \forall n\neq 0,\lambda^n\neq 1$. So $\phi_\lambda$ is not faithful if and only if $\lambda=\e^{q\pi}$ for some $q\in \Q$. 
\end{solution}

\begin{problem}
    Let $V=\mathbb{R}[x]$. $\forall a \in \mathbb{R}$, let:
$$
\begin{aligned}
& L_a(f(x)):=f(a x), \quad \forall f(x) \in \mathbb{R}[x], \\
& S_a(f(x)):=f\left(\mathrm{e}^a x\right), \quad \forall f(x) \in \mathbb{R}[x] .
\end{aligned}
$$
and:
$$
\begin{array}{ll}
\varphi(a)=L_a, & \forall a \in \mathbb{R}, \\
\psi(a)=S_a, & \forall a \in \mathbb{R} .
\end{array}
$$
Question: Is $\varphi$ and $\psi$ is infinitely dimensional real representation of $(\mathbb{R},+)$ ?
\end{problem}

\begin{solution}
    $\phi(a+b)(f)(x)=f((a+b)x),\phi(a)\phi(b)(f)(x)=\phi(a)(f(bx))=f(abx)$, let $f(x)=x,a=b=1$ we get $\phi(a+b)(f)(x)=f((a+b)x)\neq f(abx)$, so $\phi$ is not representation of $\R$. 

    $\psi(a+b)(f)(x)=f(\e^{a+b}x),\psi(a)\psi(b)(f)(x)=\psi(a)(f(\e^b x))=f(\e^a \e^b x)=\psi(a+b)(f)(x)$, so $\psi$ is a representation of $(\R,+)$. 
    Obviously $\R[x]$ is infinite-dimension, so $\psi$ is infinitely dimensional. 
\end{solution}

\begin{problem}
    Give the matrix representation of $G=\<a\>$ given by the regular representation, $\rho$, of field $K$, where $\rank(a)=4$. 
\end{problem}
\begin{solution}
    $\{e,a,a^2,a^3\}$ is basis of $K[G]$. And:
    \begin{equation}
        \rho(a)e=a,\rho(a)a=a^2,\rho(a)a^2=a^3,\rho(a)a^3=e
    \end{equation}
    So matrix of $\rho(a)$ is:
    \begin{equation}
        P(a)=\left(\begin{array}{cccc}
            0 & 0 & 0 & 1\\
            1 & 0 & 0 & 0\\
            0 & 1 & 0 & 0\\
            0 & 0 & 1 & 0\\
        \end{array}\right)
    \end{equation}
    And thus $P(a^n)=P(a)^n$. Then $P$ is the matrix representation of $G$. 
\end{solution}



\end{document}
